\documentclass[11pt,a4paper]{article}

% --- Encoding + fonts ---------------------------------------------
\usepackage[T1]{fontenc}
\usepackage[utf8]{inputenc}
\usepackage{lmodern}
\usepackage{microtype}
\setlength{\emergencystretch}{3em}

% --- Math ---------------------------------------------------------
\usepackage{amsmath, amssymb, amsthm, amsfonts}

% --- Layout -------------------------------------------------------
\usepackage[a4paper,margin=2.5cm]{geometry}
\usepackage{booktabs}
\usepackage{multirow}
\usepackage{array}
\usepackage[hidelinks]{hyperref}
\usepackage[safe]{tipa}

% --- Theorems -----------------------------------------------------
\theoremstyle{plain}
\newtheorem{theorem}{Theorem}[section]
\newtheorem{proposition}[theorem]{Proposition}

\theoremstyle{definition}
\newtheorem{definition}[theorem]{Definition}
\newtheorem{example}[theorem]{Example}

\theoremstyle{remark}
\newtheorem{remark}[theorem]{Remark}

% --- Notation -----------------------------------------------------
\newcommand{\E}{\mathbb{E}}
\newcommand{\PP}{\mathbb{P}}
\newcommand{\R}{\mathbb{R}}
\newcommand{\N}{\mathbb{N}}
\newcommand{\eqdef}{\mathrel{\stackrel{\mathrm{def}}{=}}}
\newcommand{\T}{\mathsf{T}}

\title{MORIE: A Multi-Domain Scientific Computing Toolkit
       for Observational Inference, with Sociolegal,
       Signal-Processing, Cryptographic, and Spatial-Statistics
       Modules}

\author{Vansh Singh Ruhela\thanks{%
  Correspondence: \texttt{hadesllm@proton.me}.
  Code: \url{https://github.com/hadesllm/morie}
  (open-source, GPL-2.0-only).}}

\date{2026-05-09}

\begin{document}

\maketitle

\begin{abstract}
This study describes \emph{MORIE} (Methods for Observational Inference and
Robust Analysis of Interventions in Scientific Experimentation), a
dual-language (Python and R) open-source scientific-computing
toolkit released alongside this paper. MORIE is multi-domain by
construction: a single package surface supports general-purpose
causal inference and survey-weighted estimation, signal processing
and spectral analysis (with applications to forensic audio and
biomedical signals), classical and modern cryptographic
primitives, spatial statistics, statistical-physics models of
crime, classical-test-theory and item-response-theory
psychometrics, automated scrapers for Canadian carceral and police-
oversight corpora, and a primary implementation of the MRM
(McNamara-Ruhela-Medina) framework for sociolegal data
analysis~\cite{Ruhela2026MRM}. The package ships 41 built-in
datasets in a portable SQLite layer, a 36{,}476-file individual-
function namespace, a 10-screen Textual terminal user interface, a
multi-provider local-LLM chain with a vendored TurboQuant
KV-cache compression~\cite{Zandieh2026TurboQuant}, and a pure-Python
emissions tracker covering 213 countries. The current Python test
suite contains an extensive baseline test suite plus $55$ tests
covering the federal/provincial Mandela cross-jurisdiction modules,
all passing under continuous integration; the Sphinx documentation
site builds without warnings under strict mode.
\end{abstract}

\section{Introduction}
\label{sec:intro}

Scientific computing tooling is highly fragmented: causal inference,
survey statistics, spatial analysis, signal processing,
cryptography, and statistical physics each have their own ecosystems
with distinct interfaces, data conventions, and packaging norms.
Researchers who work \emph{across} these domains --- a population
that includes anyone studying carceral / criminological data
(which mixes survey design, causal estimands, spatial statistics,
and self-exciting point processes) --- pay a recurring tax in glue
code, format conversions, and re-implementation of standard
diagnostics that already exist in some other package.

MORIE addresses this gap by providing a single namespace, a
single dataset registry, and a single set of result containers
(\texttt{RichResult}) that span all of the above domains. It is
deliberately broader than its named primary application (the MRM
framework on Canadian sociolegal data) and is intended to be
useful for anyone working at the intersections of causal
inference, signal processing, cryptography, and spatial /
statistical-physics modelling.

The package is open source under the GNU General Public License
v2 (Linus copyleft, deliberately chosen over GPL-3.0 for
compatibility with the broader Linux-kernel-style ecosystem), and
is released via PyPI, CRAN, and r-universe simultaneously. The R
package mirrors the Python package function-for-function via the
\texttt{morie} R package whose source lives at
\texttt{r-package/morie/} in the same repository.

This paper documents the package's scope and architecture; the
mathematical foundations of the MRM primary application are
developed in detail in the companion paper~\cite{Ruhela2026MRM}.

\section{Architecture}
\label{sec:architecture}

\subsection{Repository layout}

The MORIE source repository follows the standard scientific-Python
layout. The Python package lives at \texttt{src/morie/} (flat
module layout, $\approx 60$ top-level modules); the R package
lives at \texttt{r-package/morie/} (standard CRAN layout with
\texttt{R/}, \texttt{man/}, \texttt{tests/}, \texttt{inst/},
\texttt{DESCRIPTION}, \texttt{NAMESPACE}). Documentation source
sits at \texttt{docs/source/}, the test suite at
\texttt{tests/}, and the continuous-integration workflows in
\texttt{.github/workflows/}. The Sphinx site is built from
\texttt{docs/source/} and published to GitHub Pages on each push
to \texttt{main}.

\subsection{The \texttt{morie.fn} namespace}

A central architectural feature is the \texttt{morie.fn}
namespace, which contains tens of thousands of individual function
files (one function per file, with shared infrastructure in
\texttt{fn/\_containers.py}, \texttt{fn/\_richresult.py},
\texttt{fn/\_registry.py}, etc.). Each file is named under a short
stable-identifier convention (typically 3--7 characters), e.g.\
\texttt{ate.py}, \texttt{difmh.py}, \texttt{kmo.py},
\texttt{huberw.py}. The intent is canonical short names per
estimator/test/kernel/weight matrix; the registry at
\texttt{fn/\_registry.py} maps short names to fully-qualified
callable paths and a one-line description used by
\texttt{morie.fn.cheatsheet(name)}. Older monolithic modules
(\texttt{causal.py}, \texttt{psymet.py}, \ldots) re-export from
\texttt{fn/} for backward compatibility, so user-facing code is
unaffected by the per-file naming.

\subsection{The \texttt{RichResult} container}

Every analysis function in MORIE returns a
\texttt{RichResult} object --- an \texttt{R}-style multi-section
result container modelled on \texttt{psych::omega} output. A
\texttt{RichResult} has structured fields for a title, summary
lines, tables, warnings, an interpretation paragraph, and a
typed payload, and inherits from \texttt{dict} so that
\texttt{isinstance(result, dict)} returns \texttt{True} for any
downstream code that expects raw dictionary output. The container
is rendered as a multi-section paragraph block by default, and
its payload is JSON-serialisable for archival.

\subsection{Result containers and the \texttt{describe()} convention}

In addition to the result, every callable that emits a
\texttt{RichResult} ships a sibling \texttt{describe\_<short>.md}
documentation file that explains the function in pedagogical
multi-section paragraphs (theory; assumptions; equations; example
output; references). The
\texttt{morie.describe(callable)} helper renders these files
inline at the REPL or in the TUI, so the user can ask, mid-session,
``what is this function actually computing?''

\subsection{Terminal user interface}

MORIE is terminal-first by design. A 10-screen TUI built on
Textual provides interactive data exploration (Datasets), pipeline
execution (Pipelines), diagnostics (Doctor), help and reference
(Help), debugging (Debug), $\approx 900$ statistical commands
(Stat), a polyglot REPL (\S\ref{sec:polyglot}), settings, and an
LLM chat (\S\ref{sec:llm}). The TUI requires no browser and runs
on a remote ssh session; it is the recommended interface for users
in secure or air-gapped environments.

\section{Data infrastructure}
\label{sec:data}

\subsection{The 41 built-in datasets}

MORIE ships 41 datasets in a portable
\texttt{morie\_datasets.db} SQLite file ($\approx 500$~MB,
tracked via Git LFS) that is bundled inside the Python package's
\texttt{data/} directory and accessible to both Python (via the
\texttt{sqlite3} stdlib module) and R (via \texttt{DBI +
RSQLite}). The dataset catalogue covers Canadian public-health
surveillance (Statistics Canada CPADS, CCS, CSADS, CSUS PUMF
microdata; Health Infobase aggregate trend tables; CIHI hospital /
health-system indicators; the 809{,}000-row indicator library), the
open Ontario Tracking Information System (OTIS) provincial
restrictive-confinement dataset, the Ontario Special Investigations
Unit canonical 65-column case CSV produced by MORIE's own
scraper~(\S\ref{sec:scrapers}), the nine Toronto Police Service
open-data crime categories, and the Statistics Canada Crime
Severity Index weights. An optional \texttt{download-bootstrap}
subcommand fetches an additional 1.5~GB of bootstrap weights
on demand.

\subsection{Dataset-agnostic estimator interfaces}

A design rule across the package is that no estimator hard-codes
a column name. Every function takes column names as keyword
arguments with sensible defaults, so the same callable runs
unchanged across CPADS, OTIS, TPS, or any user-supplied
\texttt{pandas.DataFrame}. The CPADS column contract is provided
as a canonical reference (\texttt{morie.cpads}) but is opt-in,
not enforced.

\subsection{The \texttt{DatasetRegistry}}

A central \texttt{DatasetRegistry} class maps each dataset's
short identifier (e.g.\ \texttt{otis-2025}, \texttt{cpads-2122},
\texttt{tps-assault}) to a loader function, a cached
\texttt{ColumnProfile} per column, and a suggested-role
heuristic that flags candidate \texttt{treatment / outcome /
covariate / weight / id / stratum / cluster} columns. The
registry powers the TUI's Datasets screen and the \texttt{morie
list-datasets} CLI.

\section{Causal inference and survey statistics}
\label{sec:causal}

The \texttt{causal}, \texttt{effects}, \texttt{investigation},
\texttt{ebac}, \texttt{survey}, and \texttt{inference} modules
implement a comprehensive suite of causal estimators with full
survey-weight support:

\begin{itemize}
\item Inverse-probability-weighting and H\'ajek-stabilised IPW
      (ATE, ATT, ATC).
\item Augmented IPW (doubly robust), with the standard
      semiparametric efficiency bound when both nuisance models
      are correctly specified.
\item Double / debiased machine learning (DML) with the partially
      linear (PLR) and interactive (IRM) score functions, via the
      \texttt{DoubleML} dependency~\cite{Chernozhukov2018DoubleML}.
\item Propensity-score matching and Rosenbaum--Rubin
      subclassification.
\item AIPW with a SuperLearner ensemble~\cite{vanderLaan2007SuperLearner}.
\item Instrumental variables (2SLS / partially linear IV).
\item Regression discontinuity (sharp and fuzzy).
\item Difference-in-differences, including the Bacon decomposition.
\item Sensitivity analysis: Rosenbaum bounds, the E-value, and a
      generic \texttt{verify\_chi2()} helper for replication
      auditing.
\end{itemize}

The survey-weighted machinery uses a \texttt{SurveyDesign} class
that parallels the R \texttt{survey} package
\cite{Lumley2010ComplexSurveys} and supports
stratification, clustering, probability weights, and
calibration / bootstrap variance estimation. The
Horvitz--Thompson and H\'ajek estimators are exposed as primitives
for any user wishing to roll their own design.

\section{Signal processing, cryptography, and adjacent modules}
\label{sec:adjacent}

\subsection{Signal processing and spectral analysis}

The \texttt{signal\_processing} and
\texttt{homomorphic\_deconvolution} modules expose a working stack
for forensic-audio, biomedical-signal, and seismological
applications: classical FFT and short-time Fourier transform,
multi-taper spectral estimation, wavelet decomposition (continuous
and discrete), cepstral methods, spectral subtraction,
blind-deconvolution by cepstral liftering and homomorphic
filtering, and time--frequency reassignment. The
homomorphic-deconvolution path is exposed as a separate module
because it has independent applications outside signal processing
(notably in genomics).

\subsection{Cryptography}

The \texttt{crypto} module collects classical and modern
cryptographic primitives suitable for teaching and for low-stakes
research workflows: stream ciphers, block ciphers (DES, AES, 3DES),
hash families (MD5, SHA-2, SHA-3, BLAKE2/3), HMAC, classical
substitution and transposition ciphers (Caesar, Vigen\`ere, Hill,
playfair), elliptic-curve primitives (X25519, Ed25519), and
constant-time comparison helpers. The module is not intended as a
substitute for production-grade libraries but as a teaching layer
that shares result containers with the rest of MORIE.

\subsection{Spatial statistics}

The \texttt{spatial} module exposes Moran's $I$ (global, local,
bivariate), Ripley's $K$ and $L$, the Getis--Ord $G^{\ast}$
statistic, kernel density estimation, DBSCAN clustering on
geo-coordinates, and a polygon-based Moran's $I$ implementation
for neighbourhood-aggregate data (used heavily by the TPS-
aggregated workloads in \S\ref{sec:mrm}).

\subsection{Statistical physics of crime}

A \texttt{tps\_statphysics} module implements the
Short--Brantingham--D'Orsogna reaction--diffusion crime
model~\cite{ShortBrantingham2008}, the L\'evy-flight tail-index
estimator (Hill estimator on inter-event spatial displacements),
the Bettencourt--Lobo--West urban-scaling
exponent~\cite{Bettencourt2007} fit by HC3-robust OLS, and a
Lotka--Volterra fit for paired police--crime annual time series.
A separate \texttt{tps\_hawkes\_advanced} module fits non-stationary
Hawkes self-exciting processes with sinusoidal or constant
baselines and exponential, gamma, Weibull, or Lomax (power-law)
excitation kernels, with AIC- and time-rescaling-residual-based
kernel selection~\cite{Daley2003PointProcess}; the full
likelihood-derivation, asymptotic theory after
Kwan--Chen--Dunsmuir, and Toronto Police Service application are
developed in detail in the companion paper~\cite{Ruhela2026Hawkes}.

\subsection{Psychometrics}

The \texttt{psymet} module covers Cronbach's $\alpha$, McDonald's
$\omega$, KMO sampling adequacy, Bartlett's sphericity, parallel
analysis for factor retention, composite reliability, AVE,
classical test theory item analysis, and item-response-theory
estimation for one-, two-, and three-parameter logistic models
(via \texttt{morie.fn.irt1p}, \texttt{irt2p}, \texttt{irt3p}),
with differential-item-functioning diagnostics (Mantel--Haenszel,
logistic regression).

\section{The MRM primary application}
\label{sec:mrm}

The MRM (McNamara-Ruhela-Medina) framework, developed in detail in
the companion paper~\cite{Ruhela2026MRM}, is MORIE's named
primary application. It is a multi-source framework that brings
five distinct Canadian sociolegal data streams under a coordinated
set of estimators:

\begin{enumerate}
\item Provincial Ontario Tracking Information System (OTIS)
      restrictive-confinement microdata.
\item Federal Structured Intervention Unit Implementation
      Advisory Panel reports plus the four Sprott--Doob--Iftene
      academic replications.
\item Ontario Special Investigations Unit (police-oversight)
      director's-report corpus, ingested by an automated scraper
      bundled with MORIE (\S\ref{sec:scrapers}).
\item Toronto Police Service open-data crime categories with the
      Statistics Canada Crime Severity Index.
\item Federal Corrections and Conditional Release Statistical
      Overview tables introduced into the public record by Doob's
      \emph{T-539-20} affidavit.
\end{enumerate}

MRM is implemented in MORIE as a set of \emph{MRM modules}
spanning an aggregate-IRR family, a per-individual ten-estimator
ensemble, a GEE clustering grid, a Doob $\chi^{2}$ family for
contingency tables, a Goffmanian institutional-churn analysis
suite~\cite{Goffman1961Asylums}, and a Mandela~Rules
classifier~\cite{UN2015Mandela} that operates at both the federal
and provincial levels. The framework's headline empirical finding
--- that the provincial Ontario Mandela ``torture''-classified
proportion exceeds the federal SIU rate found by Sprott and
Doob~\cite{SprottDoob2021Torture} --- is documented in
\cite{Ruhela2026MRM}.

This paper does not duplicate the mathematical content of the
MRM paper. The reader interested in the framework's theorems,
estimator derivations, or empirical replication of
Sprott--Doob--Iftene's published $\chi^{2}$ statistics should
consult~\cite{Ruhela2026MRM}.

\section{Automated scrapers}
\label{sec:scrapers}

MORIE bundles two automated data-collection pipelines, both
implemented as standalone modules with their own normalisation
schemas:

\subsection{Ontario SIU scraper}

The \texttt{morie.siu} package ingests the Ontario Special
Investigations Unit director's-report corpus end-to-end:
\begin{equation}
  \texttt{scrape\_drid} \to \texttt{parse\_html}
   \to \texttt{normalize} \to \texttt{schema}
   \to \texttt{write\_csv} \to \texttt{analyze}.
\end{equation}
The output is a 65-column canonical CSV keyed on the SIU case
number ($\mathrm{YY}$-$\mathrm{XXX}$-$\mathrm{NNN}$). Seven
\texttt{RichResult}-emitting analysers cover demographics, decision
timing, mental-health-and-race indicators, by-police-service
breakdowns, by-year breakdowns, and the case-count panel.
The pipeline is conservative on the network side ($\mathrm{drid}$
ranges scraped politely with a concurrency limit and an HTTP cache),
and all analytical outputs are reproducible from the canonical CSV.
Note that this is the \emph{Ontario police-oversight} SIU and is
distinct from the federal Structured Intervention Unit
(\S\ref{sec:mrm}); I have deliberately chosen module names
(\texttt{morie.siu} for Ontario police-oversight,
\texttt{morie.siuiap} for the federal panel) that do not begin
with a common prefix in order to make the distinction visible at
import time.

\subsection{Federal SIU IAP indexer}

The \texttt{morie.siuiap} module indexes (rather than scrapes)
the federally-appointed SIU Implementation Advisory Panel
reports, the four Sprott--Doob--Iftene CRIMSL / Schulich Law
academic publications, and Doob's $T$-$539$-$20$ Federal Court
affidavit. The module exposes citation strings, panel-member
metadata, and pointers to the source PDFs; analyses that consume
this content live in \texttt{morie.sprott\_doob} and
\texttt{morie.doob\_trends}.

\section{Polyglot REPL and language interop}
\label{sec:polyglot}

The TUI's REPL screen supports a polyglot mode that auto-detects
the language of each line (Python, R, Bash / Zsh / POSIX
\texttt{sh}) and bridges variables bidirectionally across
languages. Scalars, vectors, strings, booleans, and pandas
\texttt{DataFrame}s travel across the bridge transparently.
This is the recommended interface for users who need to mix
Python's machine-learning ecosystem with R's statistical-modelling
ecosystem in the same session. Language tags display as
\texttt{[P]} (Python), \texttt{[R]} (R), \texttt{[Z]} / \texttt{[B]}
/ \texttt{[S]} (Z-shell / Bash / POSIX shell).

\section{LLM integration}
\label{sec:llm}

A multi-provider local-LLM chain provides assistive analysis
without cloud dependencies:

\begin{enumerate}
\item Local Ollama (auto-detects available models, prefers a
      \texttt{perseus:e2b} expert model when present).
\item OllamaFreeAPI (free remote, no API key) via a vendored
      \texttt{morie.fam} client (pure-\texttt{httpx}, no
      \texttt{pydantic} / Rust dependencies).
\item Gemini, via \texttt{GEMINI\_API\_KEY}.
\item Generic OpenAI-compatible, via \texttt{LLM\_API\_BASE\_URL}
      and \texttt{LLM\_API\_KEY}.
\item A keyword-matched local fallback that runs offline.
\end{enumerate}

The \texttt{morie.quant} module implements a vendored
TurboQuant~\cite{Zandieh2026TurboQuant} KV-cache compression that
preserves the unbiased inner-product property critical for
downstream causal-inference applications, with a C reference
implementation (\texttt{quant\_ggml.c}) bridged by ctypes.
TurboQuant is data-oblivious (no calibration data) and unbiased,
which is the property that lets us use it inside the causal-
inference workflow without contaminating the estimators with the
quantization error.

\section{Emissions tracking}

The \texttt{morie.emissions} module is a pure-Python re-
implementation of the codecarbon CSV format, covering 213
countries via Statistics Canada / IEA energy-mix data, US
state-level EPA eGRID data, and Canadian provincial-grid data.
The module avoids the \texttt{pydantic} / PyO3 dependency chain
of the upstream codecarbon package, which means MORIE
emissions tracking works on Python 3.15+ where pydantic-core has
not yet shipped wheels.

\section{Reproducibility and continuous integration}

\begin{itemize}
\item \textbf{Tests.} The Python test suite contains
      an extensive baseline test suite plus $55$ tests of the
      federal/provincial Mandela cross-jurisdiction modules in the
      MRM primary, all passing under continuous integration.
\item \textbf{Documentation.} A Sphinx documentation site builds
      without warnings under strict mode (\verb|-W|), including
      auto-generated API reference for both the Python and R
      packages.
\item \textbf{Distribution.} The package is released on PyPI
      (Python; Trusted Publisher OIDC, no API token), CRAN (R;
      multi-OS \texttt{R CMD check} matrix), and r-universe
      (R; nightly binary builds from GitHub source). Each
      GitHub release auto-mints a Zenodo software DOI via the
      GitHub-Zenodo integration with metadata pinned in
      \texttt{.zenodo.json}.
\item \textbf{License.} MORIE is GPL-2.0-only (Linus copyleft).
      The Ontario SIU scraper component (\S\ref{sec:scrapers}) is
      additionally subject to the Ontario open-government licence
      on the source corpus.
\end{itemize}

\section{Limitations}

\paragraph{Sociolegal scope.} The MRM primary application is
specific to Canadian sociolegal data streams. Adapting the
framework to other jurisdictions requires user-supplied data and
domain-specific operationalisations of the Mandela Rules
classifier. I document this in \cite{Ruhela2026MRM}.

\paragraph{Cryptography.} The \texttt{crypto} module is intended
as a teaching layer and as glue for low-stakes research, not a
substitute for production cryptographic libraries. Users who need
side-channel-resistant implementations or audited code should
delegate to \texttt{cryptography} or \texttt{libsodium} bindings.

\paragraph{LLM dependencies.} The local LLM chain depends on
external services (Ollama daemon, OllamaFreeAPI community servers,
or commercial Gemini / OpenAI APIs). The keyword-matched fallback
keeps the package functional in fully offline / air-gapped
environments, but it is not a substitute for a local LLM.

\paragraph{Test coverage.} The 14k-test baseline does not yet
include systematic property-based testing for the cryptographic
primitives or fuzzing for the parser-heavy scrapers; both are on
the roadmap.

\section*{Acknowledgements}

The author thanks Prof.\ Beatrice Jauregui (University of
Toronto) for the network of mentorship that made this line of
work possible, and Prof.\ Ayobami Laniyonu (University of
Toronto) for valuable lessons over the broader period of this
work that informed the author's thinking on policing and
corrections research. The MRM methodology lineage acknowledges
the federal context provided by Anthony N.\ Doob's affidavit and
the four Sprott--Doob--Iftene independent academic reports
\cite{SprottDoob2021Torture}.

\bibliographystyle{plainnat}
\begin{thebibliography}{99}

\bibitem{Bettencourt2007}
Bettencourt, L.\ M.\ A., Lobo, J., Helbing, D., K\"uhnert, C., \&
West, G.\ B.\ (2007).
\emph{Growth, innovation, scaling, and the pace of life in
cities.}
PNAS, 104(17), 7301--7306.

\bibitem{Chernozhukov2018DoubleML}
Chernozhukov, V., Chetverikov, D., Demirer, M., Duflo, E.,
Hansen, C., Newey, W., \& Robins, J.\ (2018).
\emph{Double/debiased machine learning for treatment and
structural parameters.}
The Econometrics Journal, 21(1), C1--C68.

\bibitem{Daley2003PointProcess}
Daley, D.\ J., \& Vere-Jones, D.\ (2003).
\emph{An Introduction to the Theory of Point Processes,
Volume I.} Springer.

\bibitem{Goffman1961Asylums}
Goffman, E.\ (1961).
\emph{Asylums: Essays on the Social Situation of Mental Patients
and Other Inmates.} Anchor Books.

\bibitem{Lumley2010ComplexSurveys}
Lumley, T.\ (2010).
\emph{Complex Surveys: A Guide to Analysis Using R.} Wiley.

\bibitem{Ruhela2026MRM}
Ruhela, V.\ S.\ (2026).
\emph{The MRM Framework: A Multi-Source Statistical Foundation for
Canadian Carceral, Police, and Oversight Data, Implemented as MRM
Modules in MORIE.}
Zenodo. \url{https://doi.org/10.5281/zenodo.20096075}
(companion paper).

\bibitem{Ruhela2026MORIE}
Ruhela, V.\ S.\ (2026).
\emph{MORIE: A Multi-Domain Scientific Computing Toolkit for
Observational Inference, with Sociolegal, Signal-Processing,
Cryptographic, and Spatial-Statistics Modules.}
Zenodo. \url{https://doi.org/10.5281/zenodo.20096350} (this paper).

\bibitem{Ruhela2026Software}
Ruhela, V.\ S.\ (2026).
\emph{MORIE Toolkit: Methods for Observational Inference and
Robust Analysis of Interventions in Scientific Experimentation}
[Software, v0.1.2].
Zenodo. \url{https://doi.org/10.5281/zenodo.20111233}.
PyPI: \url{https://pypi.org/project/morie/}.
r-universe: \url{https://hadesllm.r-universe.dev/morie}.
Documentation: \url{https://hadesllm.github.io/morie/}.

\bibitem{Ruhela2026Hawkes}
Ruhela, V.\ S.\ (2026).
\emph{Criminological Hawkes Process via MORIE: Markovian and
Non-Markovian Self-Exciting Point Processes for Toronto Crime.}
Zenodo. \url{https://doi.org/10.5281/zenodo.20102198}.

\bibitem{ShortBrantingham2008}
Short, M.\ B., D'Orsogna, M.\ R., Pasour, V.\ B., Tita, G.\ E.,
Brantingham, P.\ J., Bertozzi, A.\ L., \& Chayes, L.\ B.\ (2008).
\emph{A statistical model of criminal behavior.}
Mathematical Models and Methods in Applied Sciences, 18(suppl.),
1249--1267.

\bibitem{SprottDoob2021Torture}
Sprott, J.\ B., \& Doob, A.\ N.\ (2021, February).
\emph{Solitary Confinement, Torture, and Canada's Structured
Intervention Units.} Centre for Criminology \& Sociolegal Studies,
University of Toronto.

\bibitem{UN2015Mandela}
United Nations General Assembly (2015).
\emph{United Nations Standard Minimum Rules for the Treatment of
Prisoners (the Nelson Mandela Rules).} A/Res/70/175.

\bibitem{vanderLaan2007SuperLearner}
van der Laan, M.\ J., Polley, E.\ C., \& Hubbard, A.\ E.\ (2007).
\emph{Super Learner.}
Statistical Applications in Genetics and Molecular Biology, 6(1).

\bibitem{Zandieh2026TurboQuant}
Zandieh, A., Daliri, M., Hadian, M., \& Mirrokni, V.\ (2026).
\emph{TurboQuant: Online Vector Quantization with Near-optimal
Distortion Rate.} ICLR.

\end{thebibliography}

\end{document}
